\documentclass[10pt]{scrartcl}
\usepackage[mdthm,colorsec,noauthor,noasy]{evan}
\usepackage{handout}
\newcommand{\ip}[2]{\langle #1,#2\rangle}
\DeclareMathOperator{\tr}{tr}\DeclareMathOperator{\nul}{null}\DeclareMathOperator{\range}{range}\DeclareMathOperator{\spn}{span}
\newcommand{\LL}{\mathcal L}\newcommand{\PP}{\mathcal P}\newcommand{\MM}{\mathcal M}
\title{Math 110 --- Week 7}
\subtitle{Solutions}
\author{}\date{}\ihead{}
\begin{document}
\maketitle
\begin{prob}
\textbf{HW7 P1. }\par 
\idea{$A^2=-I$ and $(B-I)^2=0$ give low-degree polynomials that kill $A$ and $B$; then rule out degree $1$.}
\begin{enumerate}[(a)]
  \ii $A^2=-I$, so $m_A\mid x^2+1$. $A$ is not a multiple of $I$ ($Ae_1=e_2$ is not a multiple of $e_1$), so $\deg m_A\ne1$. Hence $m_A=x^2+1$ over $\RR$. Over $\CC$ the definition (first $k$ with $A^k\in\spn(I,\dots,A^{k-1})$) gives the same polynomial: $x^2+1=(x-i)(x+i)$.
  \ii $(B-I)^2=\begin{psmallmatrix}0&1\\0&0\end{psmallmatrix}^2=0$ and $B\ne I$, so $m_B=(x-1)^2$.
  \ii $A$ over $\RR$: $x^2+1$ has no real roots, so it is not a product of real linear factors: not diagonalizable. Over $\CC$: distinct roots $\pm i$, diagonalizable. $B$: repeated root, not diagonalizable. This matches Homework 6.
\end{enumerate}
\end{prob}

\begin{prob}
\textbf{HW7 P2. }\par 
\idea{$D$ lowers degree, so a power of it is $0$; find the first one.}
(a) $D^4=0$ on $\PP_3(\RR)$, while $D^3x^3=6\ne0$. So $m_D\mid x^4$ and is not $x^k$ for $k\le3$: $m_D=x^4$. (b) Repeated root: not diagonalizable. (Also $E(0,D)=$ constants, of dimension $1<4$.)
\end{prob}

\begin{prob}
\textbf{HW7 P3. }\par 
\begin{enumerate}[(a)]
  \ii \idea{If $a_0=0$, then $x$ divides $m_T$; since $T$ is invertible, the factor $T$ can be cancelled, leaving a smaller polynomial that kills $T$.}
  If $a_0=0$, then $m_T=xq$ with $\deg q=d-1$, and $0=m_T(T)=Tq(T)$. Multiply by $T^{-1}$: $q(T)=0$, contradicting minimality of $m_T$ (if $d=1$, $q=1$ and $q(T)=I\ne0$).
  \ii \idea{Move the constant term to one side and factor $T$ out of the rest.}
  From $m_T(T)=0$: $T\big(T^{d-1}+a_{d-1}T^{d-2}+\dots+a_1I\big)=-a_0I$. So $T^{-1}=-\frac1{a_0}\big(T^{d-1}+a_{d-1}T^{d-2}+\dots+a_1I\big)$.
  \ii \idea{Apply (b), then factor $m_T$ to read off the eigenvalues.}
  $a_0=4$, so $T^{-1}=-\frac14(T^3+T^2-T-5I)$. $m_T=(x-1)^2(x^2+3x+4)$ (check: $1$ is a root of $m_T$ and of $m_T'=4x^3+3x^2-2x-5$). The eigenvalues are $1$ and $\frac{-3\pm i\sqrt7}2$. The root $1$ is repeated, so $T$ is not diagonalizable.
\end{enumerate}
\end{prob}

\begin{prob}
\textbf{HW7 P4. }\par 
\idea{Substituting $x\mapsto x+\lambda$ (or $x\mapsto x/c$) turns polynomials that kill $T$ into polynomials that kill $T-\lambda I$ (or $cT$), and back. So the two minimal polynomials divide each other's transforms and have the same degree.}
\begin{enumerate}[(a)]
  \ii $p(x)=m_T(x+\lambda)$ is monic of degree $d=\deg m_T$ and $p(T-\lambda I)=m_T(T)=0$, so $m_{T-\lambda I}\mid p$ and $\deg m_{T-\lambda I}\le d$. Likewise $r(x)=m_{T-\lambda I}(x-\lambda)$ is monic with $r(T)=m_{T-\lambda I}(T-\lambda I)=0$, so $d\le\deg m_{T-\lambda I}$. A monic divisor of $p$ of the same degree is $p$.
  \ii $p(x)=c^dm_T(x/c)$ is monic of degree $d$ (leading coefficient $c^d\cdot c^{-d}$) and $p(cT)=c^dm_T(T)=0$, so $m_{cT}\mid p$. The same argument with $cT$ and $1/c$ gives $d\le\deg m_{cT}$. So $m_{cT}=p$.
\end{enumerate}
\end{prob}

\begin{prob}
\textbf{HW7 P5. }\par 
\idea{The roots of $m_A$ are the eigenvalues and $\deg m_A\le2$, which leaves two candidates.}
$\deg m_A\le2$ (Cayley--Hamilton) and its only root is $\lambda$, and over $\CC$ it splits, so $m_A=x-\lambda$ or $(x-\lambda)^2$. $m_A=x-\lambda$ exactly when $A-\lambda I=0$. $A$ is diagonalizable iff $m_A$ has no repeated root iff $m_A=x-\lambda$ iff $A=\lambda I$.
\end{prob}

\begin{prob}
\textbf{HW7 P6. }\par 
\idea{Triangular, so the eigenvalues are on the diagonal; list the candidate minimal polynomials and test them.}
Eigenvalues $2,3$. $E(2)=\nul\begin{psmallmatrix}0&1&0\\0&0&0\\0&0&1\end{psmallmatrix}=\spn(e_1)$, $E(3)=\nul\begin{psmallmatrix}-1&1&0\\0&-1&0\\0&0&0\end{psmallmatrix}=\spn(e_3)$. Candidates (roots exactly $2,3$, degree $\le3$): $(x-2)(x-3)$, $(x-2)^2(x-3)$, $(x-2)(x-3)^2$. $(C-2I)(C-3I)$ and $(C-2I)(C-3I)^2$ have $(1,2)$ entry $\mp1\ne0$; $(C-2I)^2=\operatorname{diag}(0,0,1)$ and $(C-2I)^2(C-3I)=0$. So $m_C=(x-2)^2(x-3)$: not diagonalizable ($\dim E(2)+\dim E(3)=2<3$).
\end{prob}

\begin{prob}
\textbf{HW7 P7. }\par 
\idea{Lagrange interpolation: a product of factors $(x-\lambda_k)$ vanishes at every $\lambda_k$ but $\lambda_j$; divide by its value at $\lambda_j$.}
$p_j(x)=\prod_{k\ne j}\frac{x-\lambda_k}{\lambda_j-\lambda_k}$. Since $m_T$ has distinct roots, $T$ is diagonalizable and $V=\bigoplus_kE(\lambda_k,T)$. For $v\in E(\lambda_k,T)$, $p_j(T)v=p_j(\lambda_k)v$, which is $v$ if $k=j$ and $0$ otherwise. So for $v=\sum_kv_k$, $p_j(T)v=v_j$: the projection onto $E(\lambda_j,T)$ along $\bigoplus_{k\ne j}E(\lambda_k,T)$.
\end{prob}

\begin{prob}
\textbf{M1. }\par 
\idea{Find a low-degree polynomial relation, then rule out its proper divisors. In (c), $B$ is block diagonal, so $m_B$ is the lcm of the blocks' minimal polynomials (F4).}
\begin{enumerate}[(a)]
  \ii Every entry of $A^2$ is $3$, so $A^2=3A$ and $m_A\mid x(x-3)$. $A\ne0$ and $A\ne3I$, so $m_A=x(x-3)$. Distinct roots: diagonalizable.
  \ii $Pe_1=e_2$, $Pe_2=e_3$, $Pe_3=e_1$, so $P^3=I$ and $m_P\mid x^3-1$. Since $e_1,Pe_1,P^2e_1$ is a basis, no nonzero polynomial of degree $\le2$ kills $e_1$, so $m_P=x^3-1=(x-1)(x^2+x+1)$. Over $\CC$ the roots $1,\omega,\bar\omega$ are distinct: diagonalizable. Over $\RR$, $x^2+x+1$ has no real root: not diagonalizable.
  \ii $B$ is block diagonal with blocks $B_1=\begin{psmallmatrix}3&1\\0&3\end{psmallmatrix}$, $(3)$, $(5)$. $B_1\ne3I$ and $(B_1-3I)^2=0$, so $m_{B_1}=(x-3)^2$; the other blocks give $x-3$ and $x-5$. So $m_B=\operatorname{lcm}\big((x-3)^2,x-3,x-5\big)=(x-3)^2(x-5)$. Repeated root: not diagonalizable.
\end{enumerate}
\end{prob}

\begin{prob}
\textbf{M2. }\par 
\begin{enumerate}[(a)]
  \ii \idea{Test the basis vectors: $x^k$ is an eigenvector.}
  $T(x^k)=kx^k$, so $1,x,\dots,x^n$ is an eigenbasis with distinct eigenvalues $0,1,\dots,n$. $T$ is diagonalizable and $m_T=x(x-1)\cdots(x-n)$.
  \ii \idea{$S-I$ is a ``difference operator'' that lowers degree, like differentiation; so it is nilpotent, and the question is the first power that kills everything.}
  $(S-I)p=p(x+1)-p(x)$: the leading terms cancel, so $\deg$ drops by at least $1$ and constants go to $0$. Hence $(S-I)^{n+1}=0$, so $m_S\mid(x-1)^{n+1}$. Also $(S-I)x^k=(x+1)^k-x^k=kx^{k-1}+(\text{lower})$, so $(S-I)^nx^n=n!\ne0$. So $m_S=(x-1)^{n+1}$, a repeated root: not diagonalizable.
\end{enumerate}
\end{prob}

\begin{prob}
\textbf{M3. }\par 
\idea{$T^3=T$ says $x^3-x$ kills $T$; factor it.}
$m_T\mid x^3-x=x(x-1)(x+1)$, so $m_T$ is a product of distinct real linear factors and $T$ is diagonalizable. Its eigenvalues are roots of $m_T$, hence in $\{-1,0,1\}$.
\end{prob}

\begin{prob}
\textbf{M4. }\par 
\idea{A monic polynomial of degree $1$ is $x-c$, and $m_T=x-c$ says $T-cI=0$.}
If $T=cI$, then $x-c$ kills $T$ and $\deg m_T\ge1$ (as $V\ne0$), so $m_T=x-c$. Conversely, if $m_T=x-c$, then $T-cI=m_T(T)=0$.
\end{prob}

\begin{prob}
\textbf{M5. }\par 
\idea{Substitute $x\mapsto1/x$ in $m_T(T)=0$ and clear denominators by multiplying by $T^{-d}$.}
$q(x)=\frac1{a_0}x^dm_T(1/x)=\frac1{a_0}\big(1+a_{d-1}x+\dots+a_0x^d\big)$ is monic of degree $d$ ($a_0\ne0$ by HW7 P3). $q(T^{-1})=\frac1{a_0}T^{-d}m_T(T)=0$, so $m_{T^{-1}}\mid q$ and $\deg m_{T^{-1}}\le d$. Applying this to $T^{-1}$ in place of $T$ gives $d\le\deg m_{T^{-1}}$. So $m_{T^{-1}}=q$.
\end{prob}

\begin{prob}
\textbf{M9. }\par 
\idea{$m_A(A)=0$ is a rule for rewriting $A^2$ in terms of $A$ and $I$; use it to reduce any power, and solve it for $A^{-1}$.}
$A$ is not scalar, so $m_A=\chi_A=x^2-(\tr A)x+\det A=x^2-5x-2$. From $A^2-5A=2I$: $A(A-5I)=2I$, so $A^{-1}=\frac12(A-5I)=\begin{psmallmatrix}-2&1\\3/2&-1/2\end{psmallmatrix}$. $A^3=A(5A+2I)=5A^2+2A=5(5A+2I)+2A=27A+10I$.
\end{prob}

\end{document}
